% =============================================================
%  第三单元  数一数与乘法
% =============================================================
\Chapter{第三关\quad 把加法压扁}
\noindent{\small\kaishu\color{tipc}（课本第三单元：数一数与乘法，P30--P42）}

\begin{blunder}
爸爸小时候家里买了三十个鸡蛋，
爸爸蹲在厨房一个一个数，数到一半被叫走，回来又从头数。
大人说：“豆角爸爸数鸡蛋要数到天黑。”
——那时候没人告诉豆角爸爸，鸡蛋盒就是一排一排的。
\end{blunder}

\begin{gamble}{赌一把：谁先数完一盒鸡蛋}
\begin{itemize}
  \item 去超市拍一盒鸡蛋的照片（或者家里找东西排成 5 行 $\times$ 2 排）。
        \textbf{小豆角可以一个一个数，爸爸不可以。}
  \item 豆角爸爸念“5、10”，报出 10。小豆角一个一个数到 10。
        \textbf{豆角爸爸赢}，但如果小豆角能说出“因为每排 5 个”，这局算小豆角赢。
  \item 赌注：\textbf{赢的人决定周末干什么}。
\end{itemize}
\end{gamble}

\begin{bonus}
\begin{itemize}
  \item \textbf{“一排一排”寻宝}：在家里找出五处“排得整整齐齐”的东西
        （药板、蛋盒、地砖、书架、爸爸整箱的啤酒），拍照，每张照片写一个乘法算式。
        攒够十张，作成一本小册子。
  \item \textbf{全家有几倍}：筷子是勺子的几倍？椅子是桌子的几倍？
        拿真东西数，比做题有效十倍。
  \item \textbf{两分钟摆图}：给一个算式（比如 $3\times4$），两分钟内用小棒摆出两种不同的图，
        摆好拍照，一周后拿出来让小豆角再讲一遍。
\end{itemize}
\end{bonus}

\begin{unitinfo}
\textbf{课本 P30--P42}\quad \textbf{5 次课，每次 20 分钟}\quad
\textbf{建议第 8--10 周} \\[4pt]
\textbf{这一关到底在干嘛}：加法写着太长了，于是发明了乘法。
乘法不是新东西，它就是\textbf{“几个几相加”的缩写}。\\
\textbf{逻辑是这么走的}：\ 一样一样数太慢 $\to$ 一组一组数 $\to$ 写成加法
$\to$ 加法太长 $\to$ 缩写乘乘法 $\to$ 横着看竖着看是一回事 $\to$
“倍”只是“几个几”的另一种说法。\\
\textbf{这一章一句口诀都不教。}谁在这儿抢跑背口诀，到第七单元一定乱。
\end{unitinfo}

\subsection{开局先聊个事：乘法是从“嫌麻烦”里长出来的}

超市里的鸡蛋盒，一排 5 个，2 排。豆角爸爸数鸡蛋会一个一个数吗？
不会。豆角爸爸会“5、10”——两下数完。

教室里座位、乐高底板、药板、地砖、巧克力块，全是这样一格一格的。
\textbf{凡是一排一排排整齐的东西，都能用乘法数。}

这也是为什么乘法诞生得那么早：古代人收粮食、数牲口，
一个一个数实在受不了，就发明了“几个几”。

\subsection{任务 1：一组一组地数}

\begin{lesson}{1}{20}
\textbf{今天要拿下的}：小豆角体会到“成组数”比“一个一个数”快，并且能把数法写成加法。

\begin{materials}
豆子或其他颗粒 15 个、小棒。
\end{materials}

\begin{enumerate}
  \item \textbf{一堆一堆数}：把 15 颗豆子随手分 3 堆，每堆 5 颗（故意散着放）。
        让小豆角数一共几颗——小豆角会一个一个数：1、2、3……15。
        问一句：“能不能快一点？”
  \item \textbf{排整齐}：把豆子排成 3 排，每排 5 颗。让小豆角 5 个 5 个数：5、10、15。
        问：“快多少？”
  \item \textbf{写下来}：“小豆角是怎么数的？把数法写出来。”$\to 5+5+5=15$。
  \item \textbf{换个方向数}：同样这些豆子，排成 5 排、每排 3 颗。
        一列一列数 $\to 3+3+3+3+3=15$。
        \textbf{关键对比}：两个算式\textbf{长得不一样}，结果都是 15。
        问：“为什么一样？”（豆子还是那 15 颗，只是看法换了）
  \item \textbf{加难度}：4 盘，每盘 6 个（课本 P31）。
        横着 $6+6+6+6=24$；竖着 $4+4+4+4+4+4=24$。
  \item \textbf{让小豆角“烦”一次}：如果 12 盘，每盘 4 个呢？让小豆角把算式写出来。
        写到第 6 个 4 就开始烦躁了。\textbf{就停在这儿}，说：
        “明天教小豆角个短一点的写法。”
        \textbf{这个“烦”，就是乘法的动机。}千万别提前告诉小豆角答案。
\end{enumerate}

\begin{saybox}
“小豆角数了几次？每次几个？”\\
“横着数的算式和竖着数的不一样，怎么结果一样？”\\
“12 盘每盘 4 个，小豆角写累了吗？有没有更短的写法？”\\
“小豆角觉得什么样的数法能写成乘法？”（引出“每组一样多”）
\end{saybox}

\begin{passbox}
\begin{itemize}
  \item 看着图能说出“每行几个、几行”，并写出对应的加法。
  \item 小豆角能说出“每组一样多才能用乘法”——这是这一章的核心。
\end{itemize}
\end{passbox}
\end{lesson}

\subsection{任务 2：认识乘法}

\begin{lesson}{2}{20}
\textbf{今天要拿下的}：认识乘号、乘数、积；能说出一个乘法算式的两种意思。

\begin{materials}
豆子 20 颗、白纸、笔。
\end{materials}

\begin{enumerate}
  \item \textbf{接上昨天的“烦”}：“12 个 4 相加”写起来太长。
        数学家偷了个懒，写成 $4\times12$。就这么简单，没什么神秘的。
  \item \textbf{认名字}（在白板上写一个大的）：
        \begin{center}
        \begin{tikzpicture}[font=\small]
          \node (expr) at (0,0) {\Large $2\quad\times\quad3\quad=\quad6$};
          \node[below=4mm of expr, xshift=-1.05cm, text=mainc] {\scriptsize 乘数};
          \node[below=4mm of expr, xshift=0.55cm, text=mainc] {\scriptsize 乘号};
          \node[below=4mm of expr, xshift=1.55cm, text=mainc] {\scriptsize 乘数};
          \node[below=4mm of expr, xshift=3.05cm, text=mainc] {\scriptsize 积};
        \end{tikzpicture}
        \end{center}
        \textbf{读作}“2 乘 3”。（课本用“乘”，不说“乘以”。）
  \item \textbf{一个算式，两种摆法}（用豆子现场摆）：
        \begin{itemize}
          \item $2\times3$ 可以是 \textbf{3 堆、每堆 2 个}（$2+2+2$）。
          \item $2\times3$ 也可以是 \textbf{2 堆、每堆 3 个}（$3+3$）。
        \end{itemize}
        \textbf{结果一样，图不一样。}两种图都让小豆角摆、都让小豆角画下来。
  \item \textbf{课本情境}（P32 儿童乐园）：3 架小飞机，每架 2 人。
        $2\times3=6$（人）或 $3\times2=6$（人），\textbf{两个都对}。
        因为横着看是“3 个 2”，竖着看是“2 个 3”。
  \item \textbf{双向练}：
        \begin{center}
        \small
        \begin{tabular}{@{}p{6.0cm}p{6.0cm}@{}}
        \toprule
        \textbf{加法} & \textbf{乘法} \\
        \midrule
        $4+4+4$ & \blank[3.4cm] \\
        $7+7+7+7+7$ & \blank[3.4cm] \\
        $6+6+6+6$ & \blank[3.4cm] \\
        — & $5\times2$ 写成加法 \\
        — & $3\times4$ 写成加法 \\
        \bottomrule
        \end{tabular}
        \end{center}
  \item \textbf{当医生}：$3+4$ 能不能写成 $3\times4$？
        \textbf{不能}——两个加数不一样。让小豆角自己说出“乘法必须每个加数都一样”。
\end{enumerate}

\begin{saybox}
“$2\times3$ 小豆角摆的是 3 堆 2 个，还是 2 堆 3 个？都对，说说为什么。”\\
“$3+4$ 能写成乘法吗？为什么不能？”\\
“$5\times2$ 里的 5 是什么？2 是什么？”（让小豆角指着图说）
\end{saybox}

\begin{pitbox}
\textbf{三个坑}：
\begin{enumerate}[label=\arabic*、]
  \item \textbf{把所有加法都改乘法}。$3+4$ 不行。
        \textbf{一句话}：“加数一样才能改，不一样就不行。”
  \item \textbf{读法错}：读成“2 乘以 3”。\textbf{统一读“2 乘 3”。}
  \item \textbf{只会写结果，说不出意思}。考试常问“$4\times7$ 表示什么”，小豆角答不上。
        \textbf{解法}：每次写乘法算式，必须配一张图或一句话。
\end{enumerate}
\end{pitbox}

\begin{passbox}
\begin{itemize}
  \item 4 个相同加数的加法能改乘法，反过来也行。
  \item $2\times3$ 的两种意思小豆角都说得出来，各画一张图。
  \item 小豆角能判断 $3+4$ 不能改乘法，并说清理由。
\end{itemize}
\end{passbox}
\end{lesson}

\subsection{任务 3：横着看，竖着看}

\begin{lesson}{3}{20}
\textbf{今天要拿下的}：一张点子图，两种圈法，两个乘法算式。

\begin{materials}
点子图（附录 B 方格纸可以自己点）、彩笔两支。
\end{materials}

\begin{enumerate}
  \item \textbf{出示点子图}（4 行 7 列，课本 P34）：
        \begin{center}
        \begin{tikzpicture}[scale=0.46]
          \foreach \x in {0,...,6}{\foreach \y in {0,...,3}{
            \fill[mainc] (\x,\y) circle (2.6pt);}}
          \draw[pitc,thick,dashed] (-0.32,1.68) rectangle (6.32,2.32);
          \node[pitc,right,font=\scriptsize] at (6.4,2) {一行 7 个};
          \draw[sayc,thick,dashed] (4.68,-0.32) rectangle (5.32,3.32);
          \node[sayc,below,font=\scriptsize] at (5,-0.4) {一列 4 个};
        \end{tikzpicture}
        \end{center}
  \item \textbf{横着数}：“每行 7 个，4 行” $\to 7+7+7+7=28$ $\to 4\times7=28$。
  \item \textbf{竖着数}：“每列 4 个，7 列” $\to 4+4+4+4+4+4+4=28$ $\to 7\times4=28$。
  \item \textbf{让小豆角画两种圈法}：同一张图画两遍，一遍横着圈、一遍竖着圈，
        各写一个算式。\textbf{这是课本 P35 的正式要求，考试常考。}
  \item \textbf{找家里的点子}：鸡蛋盒（一排 5 个、2 排）、地砖、教室座位、
        药板、巧克力。挑一处让小豆角说“每排几个、几排”，写算式。
  \item \textbf{收一句}：“同一张图，横着数写 $4\times7$，竖着数写 $7\times4$，
        得数一样。”——不用讲“乘法交换律”这五个字。
\end{enumerate}

\begin{saybox}
“这张图小豆角还能怎么数？”\\
“$4\times7$ 和 $7\times4$ 得数一样，那图一样吗？”（图一样，圈法不一样）\\
“家里什么东西是‘一排一排’的？”
\end{saybox}

\begin{passbox}
\begin{itemize}
  \item 给任意“几行几列”的点子图，小豆角能写出\textbf{两个}乘法算式。
  \item 按要求“横圈一次、竖圈一次”，两幅图不串。
\end{itemize}
\end{passbox}
\end{lesson}

\subsection{任务 4：“倍”就是几个几}

\begin{lesson}{4}{20}
\textbf{今天要拿下的}：理解“A 是 B 的 2 倍”就是“A 里有 2 个 B”。

\begin{materials}
豆子或小棒、彩笔。
\end{materials}

\begin{enumerate}
  \item \textbf{先圈}（课本 P38 快乐的动物）：猴子 3 只，鹅 6 只。
        把鹅 3 只 3 只地圈起来：
        \begin{center}
        \begin{tikzpicture}[scale=0.5]
          \foreach \x in {0,1,2}{\fill[matc] (\x,1.4) circle(4pt);}
          \node[left,font=\scriptsize] at (-0.3,1.4) {猴子};
          \foreach \x in {0,1,2,3,4,5}{\fill[mainc] (\x,0) circle(4pt);}
          \node[left,font=\scriptsize] at (-0.3,0) {鹅};
          \draw[pitc,thick,rounded corners=3pt] (-0.28,-0.28) rectangle (2.28,0.28);
          \draw[pitc,thick,rounded corners=3pt] (2.72,-0.28) rectangle (5.28,0.28);
        \end{tikzpicture}
        \end{center}
        问：“鹅能圈出几个‘3 只’？”（2 个）$\to$ 所以鹅是猴子的 \textbf{2 倍}。
  \item \textbf{用小豆角的话讲清}：“鹅有 2 个猴子那么多，就说鹅是猴子的 2 倍。”
  \item \textbf{反复练这个句式}：“（谁）是（谁）的（几）倍。”让小豆角\textbf{手指着图}说：
        \begin{itemize}
          \item 鹅 6 只、猴子 3 只 $\to$ 鹅是猴子的 \blank 倍。
          \item 小鸟 12 只、鹅 6 只 $\to$ 小鸟是鹅的 \blank 倍。（2 倍）
          \item 小兔 2 只、小鸡 6 只 $\to$ 小鸡是兔子的 \blank 倍。（3 倍）
        \end{itemize}
  \item \textbf{跟“多几只”分清楚（重点）}：猴子 3 只、鹅 6 只。
        \begin{itemize}
          \item “鹅比猴子多几只？”$6-3=3$ 只——这是\textbf{差}。
          \item “鹅是猴子的几倍？”$6$ 里面有 2 个 3——这是 \textbf{2 倍}。
        \end{itemize}
        \textbf{同一个情境，两个问题，两个答案。}让小豆角自己说清区别。
  \item \textbf{画出来}：给“$\square$ 有 4 个，$\bigcirc$ 是 $\square$ 的 3 倍”，
        让小豆角画 $\bigcirc$：画 3 组，每组 4 个，一共 12 个。
        \textbf{能画出来，才算懂。}
\end{enumerate}

\begin{saybox}
“鹅能圈出几个 3 只？圈给豆角爸爸看。”\\
“鹅比猴子多 3 只，和鹅是猴子的 2 倍，是一件事吗？”（不是）\\
“小豆角能用‘倍’说一句家里的事吗？”（比如“小豆角的铅笔是弟弟的 3 倍”）
\end{saybox}

\begin{pitbox}
\textbf{最典型的错}：小豆角把“是 2 倍”理解成“多 2 个”。\\
\textbf{解法}：让小豆角圈。\textbf{圈出几组，就是几倍。}
只要小豆角还在算“$6-3=3$，所以是 3 倍”，说明“倍”这个概念还没进来，
退回第一步重新圈，别往下讲。
\end{pitbox}

\begin{passbox}
\begin{itemize}
  \item 给两个数量（一个正好是另一个的整倍数），小豆角能圈出组数、说出倍数。
  \item 小豆角能分清“多几只”和“是几倍”。
  \item 按“A 是 B 的 3 倍”能把 A 画出来。
\end{itemize}
\end{passbox}
\end{lesson}

\subsection{任务 5：求几倍是多少，用乘法}

\begin{lesson}{5}{20}
\textbf{今天要拿下的}：知道“一份是多少”和“几倍”，用乘法求总数；能自己编题。

\begin{materials}
纸、笔、小棒。
\end{materials}

\begin{enumerate}
  \item \textbf{课本 P39 文具店}：小兔买了 2 支铅笔，花 4 元。
        小牛买的支数是小兔的 3 倍，小牛买几支？
  \item \textbf{先圈再算}：2 支当一份，画 3 份 $\to 2\times3=6$（支）。
        \textbf{一定先画图}，别直接列式。
  \item \textbf{钱也能用倍}：小兔花 4 元，小狐狸花的钱是小兔的 4 倍
        $\to 4\times4=16$（元）。
  \item \textbf{归纳一句}：“\textbf{求一个数的几倍是多少，用乘法}：一份 $\times$ 倍数 $=$ 总数。”
  \item \textbf{三种问法摆在一起（这张表要练透）}：
        \begin{center}
        \small
        \begin{tabularx}{\linewidth}{@{}p{6.2cm}p{3.2cm}X@{}}
        \toprule
        \textbf{题目} & \textbf{列式} & \textbf{为什么} \\
        \midrule
        每盒 6 个，4 盒一共几个？ & $6\times4=24$ & 4 个 6 \\
        苹果 5 个，梨是苹果的 3 倍，梨几个？ & $5\times3=15$ & 3 个 5 \\
        红花 8 朵，比黄花多 3 朵，黄花几朵？ & $8-3=5$ &
        \textbf{这是“多几”，不是“几倍”} \\
        \bottomrule
        \end{tabularx}
        \end{center}
        \textbf{第三行是陷阱}，专门用来试小豆角有没有把“倍”当“多”。
  \item \textbf{让小豆角编题}：给算式 $6\times3=18$，编两个不同的生活故事——
        一个讲“几倍”，一个讲“几个几”，讲给豆角爸爸听。\textbf{编得出来，就是真懂。}
  \item \textbf{连线游戏}（课本 P37）：左边 $3\times6$、$4\times3$、$2\times5$，
        右边“3 个 6”“4 个 3”“2 个 5”“$6+6+6$”“$3+3+3+3$”，打乱后连。
\end{enumerate}

\begin{saybox}
“小牛的铅笔要画几份？每份几支？”\\
“这题是‘多几’还是‘几倍’？小豆角从哪个字看出来的？”（抓关键词：倍／多）\\
“$6\times3$ 小豆角能编两个不一样的故事吗？”
\end{saybox}

\begin{passbox}
\begin{itemize}
  \item 上表三种问法全对，能说清“为什么用乘法／为什么用减法”。
  \item 能看图列式，也能按算式画图。
  \item 能自己编一道“几倍”的应用题。
\end{itemize}
\end{passbox}
\end{lesson}

\subsection{翻车现场（对着看）}

\begin{center}
\small
\begin{tabularx}{\linewidth}{@{}p{3.8cm}p{4.0cm}X@{}}
\toprule
\textbf{小豆角写出来的样子} & \textbf{其实是想错了什么} & \textbf{豆角爸爸回一句什么} \\
\midrule
$3+4$ 写成 $3\times4$ & 没抓“加数必须一样” & “加数一样才能改乘法。” \\
\addlinespace[2pt]
$4\times7$ 只说一种意思 & 只记算式不记图 & 每次都要“横圈一次、竖圈一次”。 \\
\addlinespace[2pt]
“是 3 倍”算成“多 3 个” & 倍和差没分开 & 让小豆角圈，圈出几组就是几倍。 \\
\addlinespace[2pt]
读成“2 乘以 3” & 记混了不同说法 & 二年级统一读“2 乘 3”。 \\
\addlinespace[2pt]
应用题看到“倍”就加 & 没先判断是乘还是减 & 先问“是‘多几’还是‘几倍’”。 \\
\addlinespace[2pt]
算式对，单位不写 & 书写习惯 & 应用题答案必须带单位、有答语。 \\
\bottomrule
\end{tabularx}
\end{center}

\subsection{抽查：随机开枪（10 题）}

\begin{enumerate}[label=\arabic*.,itemsep=4pt]
  \item $5+5+5=$ \blank[1.4cm]，写成乘法：\blank[2.6cm]。
  \item $3+3+3+3=$ \blank[1.4cm]，写成乘法：\blank[2.6cm]。
  \item $6\times2$ 表示 \blank[1.4cm] 个 \blank[1.4cm]。
  \item $7+7+7+7+7$ 写成乘法：\blank[2.6cm]。
  \item $4\times9=36$ 里，乘数是 \blank[1.2cm] 和 \blank[1.2cm]，积是 \blank[1.2cm]。
  \item $3\times5$ 读作：\blank[3cm]。
  \item 猴子 4 只，鹅是猴子的 3 倍，鹅有 \blank[1.4cm] 只。
  \item 一本书 6 元，买 4 本要 \blank[1.4cm] 元。
  \item 每行 5 个，4 行一共 \blank[1.4cm] 个。算式 \blank[2.4cm]。
  \item 红花 8 朵，比黄花多 3 朵，黄花 \blank[1.4cm] 朵。
\end{enumerate}

\begin{ansbox}
1. 15；$5\times3$（或 $3\times5$）\quad 2. 12；$3\times4$（或 $4\times3$）\quad
3. 2 个 6（或 6 个 2）\quad 4. $7\times5$（或 $5\times7$）\quad
5. 4、9；36\quad 6. 3 乘 5\quad 7. 12\quad 8. 24\quad
9. 20；$5\times4$\quad 10. 5 \\
\textbf{第 10 题错}：“倍”和“多几”还没分开，回任务 4。
\end{ansbox}

\subsection{课后野路子}

\begin{itemize}
  \item \textbf{数一切成排的东西}：鸡蛋盒、饮料罐、地砖、药板、书架上的书。
        每天挑一样，说“每排几个、几排”，写一个乘法算式。
  \item \textbf{“豆角爸爸家有几倍”}：筷子是勺子的几倍？椅子是桌子的几倍？
        拿真东西数，比做题有效。
  \item \textbf{摆两分钟}：给一个算式（比如 $3\times4$），两分钟内用小棒摆出两种不同的图。
        摆好拍张照，一周后拿出来让小豆角再讲一遍。
  \item \textbf{电梯里数数}：一层 9 级台阶，走 6 层是多少级？（这是第七单元的伏笔）
\end{itemize}
